\section{Experimental Setup}

In this section, we describe how we conducted our experiments.

\subsection{Training Setup}

% https://huggingface.co/OpenPipe/Qwen3-14B-Instruct
We prepare our training data by processing the SWE-smith~\citep{yang2025swesmithscalingdatasoftware} using the method described in \S\ref{sec:data_creation}. Our dataset consists of a total of 39K instances from 131 repositories after pre-processing. We use a Qwen3 14B~\footnote{https://huggingface.co/OpenPipe/Qwen3-14B-Instruct} model with a modified chat template that does not remove the \texttt{<think><\textbackslash think>} during tokenization. This is important to enable training on the whole trajectory by concatenating each turn and maintain as close to the policy as possible.
Our training uses a batch size of 32 with rollouts of 4. For each rollout, we source the prompt and response tokens directly from vLLM~\cite{kwon2023efficient} to avoid retokenization drift~\footnote{https://blog.vllm.ai/2025/10/22/agent-lightning.html}. Finally, the tool responses are masked. Our experiments were conducted on a single node consisting of 8x H100 GPUs for 400 steps.

\subsection{Evaluation Task}

We evaluate \texttt{\methodname} on Software Engineering tasks: SWE-Bench Lite~\citep{jimenez2024swebenchlanguagemodelsresolve}, Verified~\footnote{https://openai.com/index/introducing-swe-bench-verified/}, and Pro~\citep{deng2025swebenchproaiagents} by calculating the F1 score. We generate the ground truth the same way we generated the training data.